\documentclass[11pt]{article}
\usepackage{reusv}

\setborderthickness{0.5pt}
\setbordermargin{1cm}
\setbordercolor{black}

\begin{document}
\drawborder
\thispagestyle{firstpage}

%% ============================================================
%%  TITLE BLOCK
%% ============================================================
\slimtitle{경로 제약의 볼록 처리}
\lightertitle{경로 제약의 볼록 처리:\\
드 카스텔조 세분화와 볼록껍질 성질}

{\normalsize\textit{Research Note No.~013 $|$ bezier-orbit-transfer-scp}}\par\medskip

\def\lighterdate{March 29, 2026}
\lighterauthor{Suwon Lee$^{\dagger}$}

\lighteraddress{$^\dagger$}{Future Mobility Control Lab, Kookmin University, Seoul, Korea}
\lighteremail{suwon.lee@kookmin.ac.kr}

%% ============================================================
%%  ABSTRACT
%% ============================================================
\begin{abstract}
\cite{report005}에서 추력 크기 제약을 볼록껍질 성질과 극값 정리로 처리하였다.
본 노트에서는 이를 \textbf{위치/경로 제약}(궤도 반경 부등식)으로 확장한다.
위치 곡선은 차수 $N{+}2$의 베지어 곡선으로서 결정변수~$\mathbf{Z}$에 아핀이므로,
드 카스텔조 세분화로 구간별 부분 제어점을 추출하면 볼록껍질 성질에 의해
\textbf{구간 전체}에서의 제약 만족을 보장할 수 있다.
선형화 오차와 볼록껍질 간극의 보수성은 $O(1/K^2)$로,
$K{=}8$ 세분화가 실용적 기본값이다.
최종 제약은 $A_{\mathrm{ub}}\,\mathrm{vec}(\mathbf{Z}) \le \mathbf{b}_{\mathrm{ub}}$
형태의 선형 부등식으로, QP/SOCP에 직접 통합된다.
수치 검증은 \cite{report014}에서 보고한다.
\end{abstract}

%% ============================================================
\section{문제 설정}
%% ============================================================

\subsection{경로 제약의 필요성}

궤도전이 최적화에서 추력 크기 제약
$\|\mathbf{u}(\tau)\| \le u_{\max}$~\cite{report005}
외에, 궤도 \textbf{반경}에 대한 부등식 제약이 필요한 경우가 있다:
\begin{equation}\label{eq:path}
  r_{\min} \;\le\; \|\mathbf{r}(\tau)\| \;\le\; r_{\max},
  \qquad \forall\,\tau \in [0,1].
\end{equation}
하한 $r_{\min}$은 대기 재진입이나 지표면 충돌을 방지하고,
상한 $r_{\max}$는 방사선 벨트 진입이나 궤도 간섭을 방지한다.

\subsection{위치 곡선의 구조}

\cite{report004,report007}에서 유도한 위치 곡선의 제어점 표현은 다음과 같다.
보조변수 승강 $\mathbf{Z} = t_f \mathbf{P}_u$를 적용하면,
위치 곡선은 차수~$N{+}2$의 베지어 곡선이 된다:
\begin{equation}\label{eq:pos_curve}
  \mathbf{r}(\tau) = \sum_{i=0}^{N+2} \mathbf{q}_{r,i}\, B_i^{N+2}(\tau),
\end{equation}
여기서 $k$번째 축의 제어점은
\begin{equation}\label{eq:qr}
  q_{r,i}^{(k)} = r_0^{(k)} + t_f v_0^{(k)} \ell_i
    + [\bar{I}_N \mathbf{Z}_{:,k}]_i
    + \phi_i^{(k)},
  \quad \ell_i = \frac{i}{N{+}2},
\end{equation}
$\bar{I}_N$는 이중 적분 행렬~\cite{report004},
$\phi_i^{(k)}$는 중력 드리프트 보정항이다.
핵심 성질은 $\mathbf{q}_r$이 $\mathbf{Z}$에 대해 \textbf{아핀}이라는 것이다:
\begin{equation}\label{eq:affine}
  \mathbf{q}_r = \mathbf{q}_{\mathrm{fixed}} + \bar{I}_N \mathbf{Z},
\end{equation}
여기서 $\mathbf{q}_{\mathrm{fixed}}$는 $\mathbf{Z}$에 무관한 고정 항이다.

\subsection{노름 제약의 비볼록성}

제약~\eqref{eq:path}은 $\|\mathbf{r}\|$에 대한 것이므로 비볼록이다.
특히 하한 $\|\mathbf{r}\| \ge r_{\min}$은
구의 \emph{외부}라는 비볼록 영역을 정의한다.
따라서 직접적인 볼록 최적화가 불가하며,
\textbf{선형화}와 \textbf{세분화}를 결합하여 볼록 근사한다.

%% ============================================================
\section{드 카스텔조 세분화와 추출 행렬}
%% ============================================================

\subsection{드 카스텔조 알고리즘}

차수~$n$의 베지어 곡선 $\mathbf{P} = [P_0, \ldots, P_n]^T$를
파라미터 $t$에서 분할하면, 왼쪽 $[0,t]$과 오른쪽 $[t,1]$의
두 부분 곡선이 생성된다. 드 카스텔조 삼각표:
\begin{equation}
  b_i^{[r]} = (1-t)\,b_i^{[r-1]} + t\,b_{i+1}^{[r-1]},
  \quad b_i^{[0]} = P_i.
\end{equation}
왼쪽 부분 곡선의 제어점은 $L_r = b_0^{[r]}$, $r=0,\ldots,n$이며,
이는 원래 제어점의 \textbf{선형 변환}이다:
\begin{equation}
  \mathbf{L} = M_{\mathrm{left}}(t) \cdot \mathbf{P},
  \quad M_{\mathrm{left}} \in \mathbb{R}^{(n+1)\times(n+1)}.
\end{equation}

\subsection{구간 추출 행렬}

임의 구간 $[a,b] \subset [0,1]$의 부분 곡선 제어점을 추출하려면
두 단계 분할을 수행한다:
\begin{enumerate}
  \item $t=b$에서 분할 $\to$ 왼쪽 $[0,b]$의 제어점 $\mathbf{L} = M_{\mathrm{left}}(b)\,\mathbf{P}$
  \item $\mathbf{L}$을 $t_2 = a/b$에서 분할 $\to$ 오른쪽이 원래 $[a,b]$에 대응
\end{enumerate}
합성하면 단일 \textbf{추출 행렬} $M_s(a,b)$를 얻는다:
\begin{equation}\label{eq:extraction}
  \boxed{\mathbf{P}_{\mathrm{sub}} = M_s(a,b) \cdot \mathbf{P},
  \quad M_s \in \mathbb{R}^{(n+1)\times(n+1)}}
\end{equation}
이 변환의 핵심 성질:
\begin{itemize}
  \item \textbf{정확성}: 근사 오차 없이 원래 곡선을 정확히 재현한다.
  \item \textbf{선형성}: $M_s$는 상수 행렬이므로, $\mathbf{P}$에 대한 선형 연산이다.
  \item \textbf{수렴}: $K \to \infty$에서 부분 제어점의 볼록껍질이 곡선 자체에 수렴한다.
\end{itemize}

%% ============================================================
\section{볼록껍질 성질의 위치 곡선 적용}
%% ============================================================

베지어 곡선의 기본 성질: 곡선은 제어점의 볼록껍질 안에 있다.
구간 $[a,b]$에서의 부분 제어점 $\mathbf{q}_{\mathrm{sub}} = M_s \cdot \mathbf{q}_r$에 대해,
\begin{equation}
  \mathbf{r}(\tau) \in \mathrm{conv}\{\mathbf{q}_{\mathrm{sub},0}, \ldots, \mathbf{q}_{\mathrm{sub},N+2}\},
  \quad \forall\,\tau \in [a,b].
\end{equation}

이때 $\mathbf{q}_r$이 $\mathbf{Z}$에 아핀이고~\eqref{eq:affine},
$M_s$가 상수 행렬이므로, 부분 제어점도 $\mathbf{Z}$에 아핀이다:
\begin{equation}\label{eq:sub_affine}
  \mathbf{q}_{\mathrm{sub}} = M_s \cdot \mathbf{q}_{\mathrm{fixed}}
    + (M_s \bar{I}_N) \cdot \mathbf{Z}.
\end{equation}
$W \triangleq M_s \bar{I}_N \in \mathbb{R}^{(N+3)\times(N+1)}$로 정의하면,
각 부분 제어점은 $\mathbf{Z}$의 선형 함수가 된다.

\textbf{비교}: \cite{report005}에서는 추력 제어점 $\mathbf{P}_u$
(차수~$N$)에 직접 볼록껍질 성질을 적용하였다.
여기서는 위치 제어점 (차수~$N{+}2$)에 적용하되,
$\bar{I}_N$을 통해 $\mathbf{Z}$와 연결한다는 차이가 있다.

%% ============================================================
\section{노름 제약의 구간별 선형화}
%% ============================================================

\subsection{반공간 근사}

비볼록 노름 제약 $\|\mathbf{r}\| \le r_{\max}$를
구간 $[\tau_k, \tau_{k+1}]$의 중점에서 선형화한다.
참조 궤적에서 중점 위치를 계산하고, 단위 방향 벡터를 구성한다:
\begin{equation}
  \hat{\mathbf{r}}_{\mathrm{mid}} =
    \frac{\mathbf{r}_{\mathrm{ref}}(\tau_{\mathrm{mid}})}
         {\|\mathbf{r}_{\mathrm{ref}}(\tau_{\mathrm{mid}})\|}.
\end{equation}
선형화된 제약:
\begin{equation}\label{eq:linearized}
  \boxed{
    \hat{\mathbf{r}}_{\mathrm{mid}} \cdot \mathbf{q}_{\mathrm{sub},i}
    \;\le\; r_{\max}, \quad \forall\,i = 0,\ldots,N{+}2
  }
\end{equation}
볼록껍질 성질에 의해, 모든 부분 제어점이~\eqref{eq:linearized}를
만족하면 구간 내 \textbf{모든}~$\tau$에서
$\hat{\mathbf{r}}_{\mathrm{mid}} \cdot \mathbf{r}(\tau) \le r_{\max}$가 보장된다.

하한 $r_{\min}$에 대해서는 부등호를 반전한다:
\begin{equation}
  \hat{\mathbf{r}}_{\mathrm{mid}} \cdot \mathbf{q}_{\mathrm{sub},i}
  \;\ge\; r_{\min}, \quad \forall\,i.
\end{equation}

\subsection{보수성 분석}

선형화에는 두 가지 보수성 원인이 있다:
\begin{enumerate}
  \item \textbf{볼록껍질 간극}: 부분 제어점의 볼록껍질이
        실제 곡선보다 넓다. 세분화 구간 수~$K$ 증가 시
        간극이 $O(1/K^2)$로 감소한다.
  \item \textbf{반공간 간극}: 구 $\|\mathbf{r}\| = r_{\max}$와
        반공간 $\hat{\mathbf{r}} \cdot \mathbf{r} = r_{\max}$의 차이.
        구간이 좁을수록 ($K$ 증가) 궤도 방향 변화가 작아 간극이 줄어든다.
\end{enumerate}
두 효과 모두 $K$에 의해 제어되므로, 결합 보수성은 $O(1/K^2)$이다.

\subsection{제약 수}

구간당 $(N{+}3)$개의 부분 제어점이 있고, $K$개 구간이 있으므로,
한쪽 제약($r_{\max}$ 또는 $r_{\min}$)당 총 제약 수는
\begin{equation}
  n_{\mathrm{constr}} = K \times (N{+}3).
\end{equation}
$r_{\min}$과 $r_{\max}$ 모두 활성이면 $2K(N{+}3)$개이다.
$K{=}8$, $N{=}12$일 때 한쪽 120개, 양쪽 240개로,
QP/SOCP 솔버의 부하로는 무시할 수 있는 수준이다.

%% ============================================================
\section{최종 제약 정식화}
%% ============================================================

\subsection{행렬 조립}

구간 $s$에 대해 다음을 계산한다:
\begin{enumerate}
  \item 추출 행렬: $M_s = M_s(\tau_{s-1}, \tau_s) \in \mathbb{R}^{(N+3)\times(N+3)}$
  \item 합성 행렬: $W = M_s \bar{I}_N \in \mathbb{R}^{(N+3)\times(N+1)}$
  \item 고정 부분: $\mathbf{q}_{\mathrm{fixed,sub}} = M_s \mathbf{q}_{\mathrm{fixed}} \in \mathbb{R}^{(N+3)\times 3}$
\end{enumerate}
$i$번째 부분 제어점에 대한 제약~\eqref{eq:linearized}를
$\mathrm{vec}(\mathbf{Z})$에 대해 전개하면:
\begin{equation}
  \sum_{k=1}^{3} \hat{r}_{\mathrm{mid}}^{(k)} W_{i,:} \mathbf{Z}_{:,k}
  \;\le\; r_{\max} - \hat{\mathbf{r}}_{\mathrm{mid}} \cdot \mathbf{q}_{\mathrm{fixed,sub},i}.
\end{equation}
$\mathrm{vec}(\mathbf{Z}) = [Z_{:,1};\, Z_{:,2};\, Z_{:,3}]$로 정의하면,
$A_{\mathrm{ub}}$의 한 행은
\begin{equation}
  [\hat{r}_1 W_{i,:},\; \hat{r}_2 W_{i,:},\; \hat{r}_3 W_{i,:}]
  \in \mathbb{R}^{3(N+1)},
\end{equation}
$b_{\mathrm{ub}}$의 대응 성분은
$r_{\max} - \hat{\mathbf{r}}_{\mathrm{mid}} \cdot \mathbf{q}_{\mathrm{fixed,sub},i}$이다.

최종 형태:
\begin{equation}\label{eq:final}
  \boxed{A_{\mathrm{ub}}\,\mathrm{vec}(\mathbf{Z}) \;\le\; \mathbf{b}_{\mathrm{ub}}}
\end{equation}
이는 기존 QP/SOCP 부분문제~\cite{report006}에 선형 부등식 제약으로 직접 추가된다.

\subsection{SCP 반복과의 통합}

경로 제약은 SCP 반복의 3번째 반복부터 활성화한다.
처음 2회는 참조 궤적이 불안정하여 선형화 방향 $\hat{\mathbf{r}}_{\mathrm{mid}}$가
신뢰할 수 없기 때문이다.
매 반복마다 현재 참조 궤적에서 $\hat{\mathbf{r}}_{\mathrm{mid}}$를 갱신하므로,
SCP가 수렴하면 선형화 방향도 수렴한다.

%% ============================================================
\section{요약}
%% ============================================================

본 노트의 핵심 기여를 정리한다:
\begin{itemize}
  \item 드 카스텔조 세분화는 추출 행렬~$M_s$를 통해
        부분 제어점을 원래 제어점의 \textbf{선형 변환}으로 제공한다.
  \item 위치 곡선의 부분 제어점에 볼록껍질 성질을 적용하면,
        격자점 평가 없이 구간 \textbf{전체}에서 제약 만족을 보장한다.
  \item 중점 선형화로 비볼록 노름 제약을 반공간 제약으로 변환하며,
        결합 보수성은 $O(1/K^2)$이다.
  \item 최종 제약~\eqref{eq:final}은 $\mathrm{vec}(\mathbf{Z})$에 대한
        선형 부등식으로, QP/SOCP 프레임워크에 직접 통합된다.
  \item \cite{report005}의 추력 볼록껍질 제약과 상보적 관계로,
        추력 크기와 궤도 반경을 모두 볼록 프레임워크 안에서 처리한다.
\end{itemize}

%% ============================================================
\bibliographystyle{IEEEtran}
\bibliography{references}

\end{document}
